%%%%%%%%%%%%%%%%%%%%%%%%%%%%%%%%%%%%%%%%%%%%%%%%%%%%%%%%%%%%
% 5.12 OPERATOR THEORY
% The Composition of Advanced Mathematics
%%%%%%%%%%%%%%%%%%%%%%%%%%%%%%%%%%%%%%%%%%%%%%%%%%%%%%%%%%%%

\section{Operator Theory}
\label{sec:operator-theory}

\prerequisite{
Linear algebra, real and complex analysis, metric spaces, normed spaces,
Banach spaces, Hilbert spaces, bounded linear operators, dual spaces,
compactness, and basic functional analysis.
}

\learningobjectives{
By the end of this section, the reader should be able to:
\begin{itemize}
    \item distinguish bounded and unbounded linear operators;
    \item compute and interpret operator norms;
    \item understand kernels, ranges, inverses, and compositions;
    \item work with operator algebras;
    \item define the spectrum and resolvent of an operator;
    \item distinguish eigenvalues from the full spectrum;
    \item classify point, continuous, and residual spectrum;
    \item use the spectral radius and spectral radius formula;
    \item understand adjoint, self-adjoint, normal, unitary, and positive operators;
    \item apply the spectral theorem in finite and infinite dimensions;
    \item understand projections and orthogonal decompositions;
    \item work with compact operators and their spectral properties;
    \item understand Fredholm operators and the Fredholm alternative;
    \item recognize shift, multiplication, differential, and integral operators;
    \item understand closed and densely defined operators;
    \item recognize the role of operator theory in differential equations,
          quantum mechanics, numerical analysis, and spectral theory.
\end{itemize}
}

%%%%%%%%%%%%%%%%%%%%%%%%%%%%%%%%%%%%%%%%%%%%%%%%%%%%%%%%%%%%
\subsection{What Is Operator Theory?}
\label{subsec:what-is-operator-theory}
%%%%%%%%%%%%%%%%%%%%%%%%%%%%%%%%%%%%%%%%%%%%%%%%%%%%%%%%%%%%

Linear algebra studies linear maps between finite-dimensional vector
spaces.

Operator theory extends this study to spaces of functions and sequences,
often of infinite dimension.

The central object is

\[
\boxed{
T:X\to Y,
}
\]

where \(X\) and \(Y\) are normed, Banach, or Hilbert spaces.

The map \(T\) is usually linear:

\[
\boxed{
T(\alpha x+\beta y)
=
\alpha Tx+\beta Ty.
}
\]

Operator theory investigates the algebraic, analytic, and geometric
properties of such maps.

%%%%%%%%%%%%%%%%%%%%%%%%%%%%%%%%%%%%%%%%%%%%%%%%%%%%%%%%%%%%
\subsection{Operators Versus Matrices}
\label{subsec:operators-versus-matrices}
%%%%%%%%%%%%%%%%%%%%%%%%%%%%%%%%%%%%%%%%%%%%%%%%%%%%%%%%%%%%

A matrix

\[
A\in\mathbb{C}^{m\times n}
\]

defines a linear transformation

\[
A:\mathbb{C}^n\to\mathbb{C}^m.
\]

In infinite-dimensional spaces, operators play the role of matrices.

Examples include

\[
\boxed{
Tf=f',
}
\]

\[
\boxed{
Tf(x)=x f(x),
}
\]

and

\[
\boxed{
Tf(x)
=
\int K(x,y)f(y)\,dy.
}
\]

Thus differentiation, multiplication, integration, and many differential
equations can be studied as operator problems.

%%%%%%%%%%%%%%%%%%%%%%%%%%%%%%%%%%%%%%%%%%%%%%%%%%%%%%%%%%%%
\subsection{Bounded Linear Operators}
\label{subsec:bounded-linear-operators}
%%%%%%%%%%%%%%%%%%%%%%%%%%%%%%%%%%%%%%%%%%%%%%%%%%%%%%%%%%%%

\begin{definitionbox}{Bounded Linear Operator}
Let \(X\) and \(Y\) be normed spaces. A linear operator

\[
T:X\to Y
\]

is bounded if there exists \(C\geq0\) such that

\[
\boxed{
\|Tx\|_Y
\leq
C\|x\|_X
}
\]

for every \(x\in X\).
\end{definitionbox}

The set of bounded linear operators from \(X\) to \(Y\) is denoted

\[
\boxed{
\mathcal{B}(X,Y).
}
\]

When \(X=Y\),

\[
\boxed{
\mathcal{B}(X)
=
\mathcal{B}(X,X).
}
\]

%%%%%%%%%%%%%%%%%%%%%%%%%%%%%%%%%%%%%%%%%%%%%%%%%%%%%%%%%%%%
\subsection{Operator Norm}
\label{subsec:operator-norm}
%%%%%%%%%%%%%%%%%%%%%%%%%%%%%%%%%%%%%%%%%%%%%%%%%%%%%%%%%%%%

For \(T\in\mathcal{B}(X,Y)\), define

\[
\boxed{
\|T\|
=
\sup_{\|x\|_X\leq1}
\|Tx\|_Y.
}
\]

Equivalently,

\[
\boxed{
\|T\|
=
\sup_{\|x\|_X=1}
\|Tx\|_Y
}
\]

when \(X\neq\{0\}\).

The operator norm is the smallest constant \(C\) satisfying

\[
\boxed{
\|Tx\|
\leq
C\|x\|.
}
\]

Therefore,

\[
\boxed{
\|Tx\|
\leq
\|T\|\|x\|.
}
\]

%%%%%%%%%%%%%%%%%%%%%%%%%%%%%%%%%%%%%%%%%%%%%%%%%%%%%%%%%%%%
\subsection{Boundedness and Continuity}
\label{subsec:boundedness-continuity-operators}
%%%%%%%%%%%%%%%%%%%%%%%%%%%%%%%%%%%%%%%%%%%%%%%%%%%%%%%%%%%%

For linear operators, boundedness and continuity are equivalent.

\begin{theorem}
Let \(T:X\to Y\) be linear between normed spaces. The following are
equivalent:
\begin{enumerate}
    \item \(T\) is bounded;
    \item \(T\) is continuous;
    \item \(T\) is continuous at \(0\).
\end{enumerate}
\end{theorem}

This equivalence is special to linear maps.

%%%%%%%%%%%%%%%%%%%%%%%%%%%%%%%%%%%%%%%%%%%%%%%%%%%%%%%%%%%%
\subsection{Algebra of Bounded Operators}
\label{subsec:operator-algebra}
%%%%%%%%%%%%%%%%%%%%%%%%%%%%%%%%%%%%%%%%%%%%%%%%%%%%%%%%%%%%

If

\[
S,T\in\mathcal{B}(X),
\]

then

\[
S+T,
\qquad
\alpha T,
\qquad
ST
\]

are bounded operators.

Moreover,

\[
\boxed{
\|S+T\|
\leq
\|S\|+\|T\|,
}
\]

and

\[
\boxed{
\|ST\|
\leq
\|S\|\|T\|.
}
\]

Thus \(\mathcal{B}(X)\) forms an algebra under operator addition,
scalar multiplication, and composition.

If \(X\) is Banach, then

\[
\boxed{
\mathcal{B}(X)
}
\]

is itself a Banach algebra.

%%%%%%%%%%%%%%%%%%%%%%%%%%%%%%%%%%%%%%%%%%%%%%%%%%%%%%%%%%%%
\subsection{Identity and Zero Operators}
\label{subsec:identity-zero-operators}
%%%%%%%%%%%%%%%%%%%%%%%%%%%%%%%%%%%%%%%%%%%%%%%%%%%%%%%%%%%%

The identity operator is

\[
\boxed{
Ix=x.
}
\]

It satisfies

\[
\boxed{
IT=TI=T.
}
\]

If \(X\neq\{0\}\),

\[
\boxed{
\|I\|=1.
}
\]

The zero operator is

\[
\boxed{
0x=0.
}
\]

These are the multiplicative and additive identities of the operator
algebra.

%%%%%%%%%%%%%%%%%%%%%%%%%%%%%%%%%%%%%%%%%%%%%%%%%%%%%%%%%%%%
\subsection{Kernel and Range}
\label{subsec:operator-kernel-range}
%%%%%%%%%%%%%%%%%%%%%%%%%%%%%%%%%%%%%%%%%%%%%%%%%%%%%%%%%%%%

For a linear operator

\[
T:X\to Y,
\]

the kernel is

\[
\boxed{
\ker T
=
\{x\in X:Tx=0\}.
}
\]

The range is

\[
\boxed{
\operatorname{Ran}(T)
=
\{Tx:x\in X\}.
}
\]

The kernel measures failure of injectivity.

The range measures which outputs can be produced.

Thus

\[
\boxed{
T\text{ injective}
\iff
\ker T=\{0\},
}
\]

and

\[
\boxed{
T\text{ surjective}
\iff
\operatorname{Ran}(T)=Y.
}
\]

%%%%%%%%%%%%%%%%%%%%%%%%%%%%%%%%%%%%%%%%%%%%%%%%%%%%%%%%%%%%
\subsection{Invertible Operators}
\label{subsec:invertible-operators}
%%%%%%%%%%%%%%%%%%%%%%%%%%%%%%%%%%%%%%%%%%%%%%%%%%%%%%%%%%%%

An operator

\[
T:X\to X
\]

is invertible if there exists

\[
T^{-1}:X\to X
\]

such that

\[
\boxed{
T^{-1}T
=
TT^{-1}
=
I.
}
\]

For bounded operators on Banach spaces, the inverse mapping theorem
implies an important fact.

\begin{theorem}[Bounded Inverse Theorem]
If \(X\) and \(Y\) are Banach spaces and

\[
T:X\to Y
\]

is bounded, linear, and bijective, then

\[
\boxed{
T^{-1}:Y\to X
}
\]

is bounded.
\end{theorem}

%%%%%%%%%%%%%%%%%%%%%%%%%%%%%%%%%%%%%%%%%%%%%%%%%%%%%%%%%%%%
\subsection{Neumann Series}
\label{subsec:neumann-series}
%%%%%%%%%%%%%%%%%%%%%%%%%%%%%%%%%%%%%%%%%%%%%%%%%%%%%%%%%%%%

Suppose

\[
T\in\mathcal{B}(X)
\]

and

\[
\boxed{
\|T\|<1.
}
\]

Then

\[
I-T
\]

is invertible and

\[
\boxed{
(I-T)^{-1}
=
\sum_{n=0}^{\infty}T^n.
}
\]

This is the operator version of the geometric series

\[
\frac{1}{1-z}
=
\sum_{n=0}^{\infty}z^n.
\]

The series converges in operator norm.

%%%%%%%%%%%%%%%%%%%%%%%%%%%%%%%%%%%%%%%%%%%%%%%%%%%%%%%%%%%%
\subsection{Worked Example: Inverting a Small Perturbation}
\label{subsec:worked-neumann}
%%%%%%%%%%%%%%%%%%%%%%%%%%%%%%%%%%%%%%%%%%%%%%%%%%%%%%%%%%%%

\begin{workedexamplebox}{Neumann-Series Inverse}

Suppose

\[
T\in\mathcal{B}(X)
\]

and

\[
\|T\|=\frac12.
\]

Since

\[
\|T\|<1,
\]

the operator \(I-T\) is invertible.

Its inverse is

\[
(I-T)^{-1}
=
I+T+T^2+T^3+\cdots.
\]

Furthermore,

\[
\|(I-T)^{-1}\|
\leq
\sum_{n=0}^{\infty}\|T\|^n.
\]

Therefore,

\begin{answerbox}
\[
\boxed{
(I-T)^{-1}
=
\sum_{n=0}^{\infty}T^n,
\qquad
\|(I-T)^{-1}\|
\leq2.
}
\]
\end{answerbox}
\end{workedexamplebox}

%%%%%%%%%%%%%%%%%%%%%%%%%%%%%%%%%%%%%%%%%%%%%%%%%%%%%%%%%%%%
\subsection{Eigenvalues and Eigenvectors}
\label{subsec:operator-eigenvalues}
%%%%%%%%%%%%%%%%%%%%%%%%%%%%%%%%%%%%%%%%%%%%%%%%%%%%%%%%%%%%

A nonzero vector \(x\in X\) is an eigenvector of \(T\) with eigenvalue
\(\lambda\) if

\[
\boxed{
Tx=\lambda x.
}
\]

Equivalently,

\[
\boxed{
(T-\lambda I)x=0.
}
\]

Therefore \(\lambda\) is an eigenvalue precisely when

\[
\boxed{
\ker(T-\lambda I)\neq\{0\}.
}
\]

In infinite dimensions, eigenvalues alone do not capture all important
spectral behavior.

%%%%%%%%%%%%%%%%%%%%%%%%%%%%%%%%%%%%%%%%%%%%%%%%%%%%%%%%%%%%
\subsection{Resolvent Set}
\label{subsec:resolvent-set}
%%%%%%%%%%%%%%%%%%%%%%%%%%%%%%%%%%%%%%%%%%%%%%%%%%%%%%%%%%%%

Let

\[
T\in\mathcal{B}(X)
\]

where \(X\) is a complex Banach space.

\begin{definitionbox}{Resolvent Set}
The resolvent set of \(T\) is

\[
\boxed{
\rho(T)
=
\{
\lambda\in\mathbb{C}:
T-\lambda I
\text{ is invertible in }\mathcal{B}(X)
\}.
}
\]
\end{definitionbox}

The lowercase Greek letter \(\rho\) is \emph{rho}.

For

\[
\lambda\in\rho(T),
\]

the resolvent operator is

\[
\boxed{
R(\lambda,T)
=
(T-\lambda I)^{-1}.
}
\]

Some texts instead use \((\lambda I-T)^{-1}\); either convention is
acceptable when used consistently.

%%%%%%%%%%%%%%%%%%%%%%%%%%%%%%%%%%%%%%%%%%%%%%%%%%%%%%%%%%%%
\subsection{Spectrum}
\label{subsec:operator-spectrum}
%%%%%%%%%%%%%%%%%%%%%%%%%%%%%%%%%%%%%%%%%%%%%%%%%%%%%%%%%%%%

\begin{definitionbox}{Spectrum}
The spectrum of \(T\) is

\[
\boxed{
\sigma(T)
=
\mathbb{C}\setminus\rho(T).
}
\]
\end{definitionbox}

The lowercase Greek letter \(\sigma\) is \emph{sigma}.

Thus

\[
\boxed{
\lambda\in\sigma(T)
}
\]

whenever

\[
T-\lambda I
\]

fails to have a bounded inverse defined on all of \(X\).

Every eigenvalue belongs to the spectrum, but the spectrum may contain
points that are not eigenvalues.

%%%%%%%%%%%%%%%%%%%%%%%%%%%%%%%%%%%%%%%%%%%%%%%%%%%%%%%%%%%%
\subsection{Spectrum Versus Eigenvalues}
\label{subsec:spectrum-versus-eigenvalues}
%%%%%%%%%%%%%%%%%%%%%%%%%%%%%%%%%%%%%%%%%%%%%%%%%%%%%%%%%%%%

In finite-dimensional linear algebra,

\[
\boxed{
\text{spectrum}
=
\text{set of eigenvalues}.
}
\]

In infinite-dimensional spaces, this may fail.

An operator may have

\[
\boxed{
\lambda\in\sigma(T)
}
\]

even though

\[
\boxed{
\ker(T-\lambda I)=\{0\}.
}
\]

The obstruction may instead be failure of surjectivity or failure of the
inverse to be bounded.

This is one of the fundamental differences between matrix theory and
infinite-dimensional operator theory.

%%%%%%%%%%%%%%%%%%%%%%%%%%%%%%%%%%%%%%%%%%%%%%%%%%%%%%%%%%%%
\subsection{Types of Spectrum}
\label{subsec:types-of-spectrum}
%%%%%%%%%%%%%%%%%%%%%%%%%%%%%%%%%%%%%%%%%%%%%%%%%%%%%%%%%%%%

For a bounded operator on a Banach space, the spectrum is commonly
divided into several parts.

The point spectrum consists of eigenvalues:

\[
\boxed{
\sigma_p(T)
=
\{
\lambda:
T-\lambda I
\text{ is not injective}
\}.
}
\]

The continuous spectrum consists, under a standard convention, of
\(\lambda\) for which \(T-\lambda I\) is injective and has dense range,
but is not surjective.

The residual spectrum consists of \(\lambda\) for which \(T-\lambda I\)
is injective but its range is not dense.

Different texts use slightly different terminology, so the convention
should always be stated.

%%%%%%%%%%%%%%%%%%%%%%%%%%%%%%%%%%%%%%%%%%%%%%%%%%%%%%%%%%%%
\subsection{Basic Spectral Theorem for Banach Algebras}
\label{subsec:basic-spectrum-properties}
%%%%%%%%%%%%%%%%%%%%%%%%%%%%%%%%%%%%%%%%%%%%%%%%%%%%%%%%%%%%

\begin{theorem}
Let \(T\in\mathcal{B}(X)\), where \(X\) is a nonzero complex Banach space.
Then

\[
\boxed{
\sigma(T)
}
\]

is nonempty and compact.

Moreover,

\[
\boxed{
\sigma(T)
\subseteq
\{
\lambda\in\mathbb{C}:|\lambda|\leq\|T\|
\}.
}
\]
\end{theorem}

Thus every bounded operator on a nonzero complex Banach space has
nonempty spectrum.

%%%%%%%%%%%%%%%%%%%%%%%%%%%%%%%%%%%%%%%%%%%%%%%%%%%%%%%%%%%%
\subsection{Why the Spectrum Lies in a Disk}
\label{subsec:spectral-disk}
%%%%%%%%%%%%%%%%%%%%%%%%%%%%%%%%%%%%%%%%%%%%%%%%%%%%%%%%%%%%

Suppose

\[
|\lambda|>\|T\|.
\]

Then

\[
T-\lambda I
=
-\lambda
\left(
I-\frac{T}{\lambda}
\right).
\]

Since

\[
\left\|
\frac{T}{\lambda}
\right\|
<1,
\]

the Neumann series shows that

\[
I-\frac{T}{\lambda}
\]

is invertible.

Therefore

\[
\lambda\notin\sigma(T).
\]

Hence

\[
\boxed{
\sigma(T)
\subseteq
\overline{B(0,\|T\|)}.
}
\]

%%%%%%%%%%%%%%%%%%%%%%%%%%%%%%%%%%%%%%%%%%%%%%%%%%%%%%%%%%%%
\subsection{Spectral Radius}
\label{subsec:spectral-radius}
%%%%%%%%%%%%%%%%%%%%%%%%%%%%%%%%%%%%%%%%%%%%%%%%%%%%%%%%%%%%

\begin{definitionbox}{Spectral Radius}
The spectral radius of \(T\) is

\[
\boxed{
r(T)
=
\sup\{
|\lambda|:\lambda\in\sigma(T)
\}.
}
\]
\end{definitionbox}

Since

\[
\sigma(T)
\subseteq
\{
|\lambda|\leq\|T\|
\},
\]

we obtain

\[
\boxed{
r(T)\leq\|T\|.
}
\]

%%%%%%%%%%%%%%%%%%%%%%%%%%%%%%%%%%%%%%%%%%%%%%%%%%%%%%%%%%%%
\subsection{Spectral Radius Formula}
\label{subsec:spectral-radius-formula}
%%%%%%%%%%%%%%%%%%%%%%%%%%%%%%%%%%%%%%%%%%%%%%%%%%%%%%%%%%%%

A fundamental result is

\[
\boxed{
r(T)
=
\lim_{n\to\infty}
\|T^n\|^{1/n}.
}
\]

Equivalently,

\[
\boxed{
r(T)
=
\inf_{n\geq1}
\|T^n\|^{1/n}.
}
\]

This formula connects the spectrum with the long-term growth of powers
of an operator.

%%%%%%%%%%%%%%%%%%%%%%%%%%%%%%%%%%%%%%%%%%%%%%%%%%%%%%%%%%%%
\subsection{Spectral Mapping}
\label{subsec:spectral-mapping}
%%%%%%%%%%%%%%%%%%%%%%%%%%%%%%%%%%%%%%%%%%%%%%%%%%%%%%%%%%%%

For a polynomial \(p\),

\[
\boxed{
\sigma(p(T))
=
p(\sigma(T)).
}
\]

This is the polynomial spectral mapping theorem.

For example,

\[
\boxed{
\sigma(T^2)
=
\{
\lambda^2:
\lambda\in\sigma(T)
\}.
}
\]

More advanced functional calculi extend this principle beyond
polynomials.

%%%%%%%%%%%%%%%%%%%%%%%%%%%%%%%%%%%%%%%%%%%%%%%%%%%%%%%%%%%%
\subsection{Functional Calculus Preview}
\label{subsec:functional-calculus-preview}
%%%%%%%%%%%%%%%%%%%%%%%%%%%%%%%%%%%%%%%%%%%%%%%%%%%%%%%%%%%%

For a polynomial

\[
p(z)
=
a_0+a_1z+\cdots+a_nz^n,
\]

define

\[
\boxed{
p(T)
=
a_0I+a_1T+\cdots+a_nT^n.
}
\]

This suggests the more general idea of defining

\[
\boxed{
f(T)
}
\]

for functions \(f\).

Examples include

\[
\boxed{
e^T
=
\sum_{n=0}^{\infty}
\frac{T^n}{n!}.
}
\]

Functional calculus is a major tool in spectral theory and differential
equations.

%%%%%%%%%%%%%%%%%%%%%%%%%%%%%%%%%%%%%%%%%%%%%%%%%%%%%%%%%%%%
\subsection{Operators on Hilbert Spaces}
\label{subsec:operators-hilbert-spaces}
%%%%%%%%%%%%%%%%%%%%%%%%%%%%%%%%%%%%%%%%%%%%%%%%%%%%%%%%%%%%

Hilbert spaces provide additional structure through the inner product

\[
\langle x,y\rangle.
\]

This allows one to define the adjoint of an operator and important
operator classes such as

\[
\boxed{
\text{self-adjoint},
\quad
\text{normal},
\quad
\text{unitary},
\quad
\text{positive}.
}
\]

These classes are infinite-dimensional analogues of familiar matrix
classes.

%%%%%%%%%%%%%%%%%%%%%%%%%%%%%%%%%%%%%%%%%%%%%%%%%%%%%%%%%%%%
\subsection{Adjoint Operator}
\label{subsec:adjoint-operator}
%%%%%%%%%%%%%%%%%%%%%%%%%%%%%%%%%%%%%%%%%%%%%%%%%%%%%%%%%%%%

Let \(H\) be a Hilbert space and

\[
T\in\mathcal{B}(H).
\]

\begin{definitionbox}{Adjoint Operator}
The adjoint \(T^*\) is the unique bounded operator satisfying

\[
\boxed{
\langle Tx,y\rangle
=
\langle x,T^*y\rangle
}
\]

for every \(x,y\in H\).
\end{definitionbox}

The notation \(T^*\) is read ``\(T\)-star.''

%%%%%%%%%%%%%%%%%%%%%%%%%%%%%%%%%%%%%%%%%%%%%%%%%%%%%%%%%%%%
\subsection{Properties of the Adjoint}
\label{subsec:adjoint-properties}
%%%%%%%%%%%%%%%%%%%%%%%%%%%%%%%%%%%%%%%%%%%%%%%%%%%%%%%%%%%%

For bounded operators \(S,T\) and scalar \(\alpha\),

\[
\boxed{
(S+T)^*
=
S^*+T^*,
}
\]

\[
\boxed{
(\alpha T)^*
=
\overline{\alpha}T^*,
}
\]

\[
\boxed{
(ST)^*
=
T^*S^*,
}
\]

and

\[
\boxed{
(T^*)^*=T.
}
\]

Also,

\[
\boxed{
\|T^*\|=\|T\|.
}
\]

A particularly important identity is

\[
\boxed{
\|T^*T\|
=
\|T\|^2.
}
\]

%%%%%%%%%%%%%%%%%%%%%%%%%%%%%%%%%%%%%%%%%%%%%%%%%%%%%%%%%%%%
\subsection{Self-Adjoint Operators}
\label{subsec:self-adjoint-operators}
%%%%%%%%%%%%%%%%%%%%%%%%%%%%%%%%%%%%%%%%%%%%%%%%%%%%%%%%%%%%

\begin{definitionbox}{Self-Adjoint Operator}
An operator \(T\in\mathcal{B}(H)\) is self-adjoint if

\[
\boxed{
T=T^*.
}
\]
\end{definitionbox}

Self-adjoint operators generalize Hermitian matrices.

For a self-adjoint operator,

\[
\boxed{
\langle Tx,x\rangle\in\R
}
\]

for every \(x\in H\).

Moreover,

\[
\boxed{
\sigma(T)\subseteq\R.
}
\]

%%%%%%%%%%%%%%%%%%%%%%%%%%%%%%%%%%%%%%%%%%%%%%%%%%%%%%%%%%%%
\subsection{Normal Operators}
\label{subsec:normal-operators}
%%%%%%%%%%%%%%%%%%%%%%%%%%%%%%%%%%%%%%%%%%%%%%%%%%%%%%%%%%%%

\begin{definitionbox}{Normal Operator}
An operator \(T\in\mathcal{B}(H)\) is normal if

\[
\boxed{
TT^*=T^*T.
}
\]
\end{definitionbox}

Important normal operators include

\[
\boxed{
\text{self-adjoint operators},
}
\]

\[
\boxed{
\text{unitary operators},
}
\]

and scalar multiples of these.

Normal operators are the natural setting for the Hilbert-space spectral
theorem.

%%%%%%%%%%%%%%%%%%%%%%%%%%%%%%%%%%%%%%%%%%%%%%%%%%%%%%%%%%%%
\subsection{Unitary Operators}
\label{subsec:unitary-operators}
%%%%%%%%%%%%%%%%%%%%%%%%%%%%%%%%%%%%%%%%%%%%%%%%%%%%%%%%%%%%

\begin{definitionbox}{Unitary Operator}
An operator \(U\in\mathcal{B}(H)\) is unitary if

\[
\boxed{
U^*U
=
UU^*
=
I.
}
\]
\end{definitionbox}

Equivalently,

\[
\boxed{
U^{-1}=U^*.
}
\]

Unitary operators preserve inner products:

\[
\boxed{
\langle Ux,Uy\rangle
=
\langle x,y\rangle.
}
\]

Therefore,

\[
\boxed{
\|Ux\|=\|x\|.
}
\]

The spectrum of a unitary operator lies on the unit circle:

\[
\boxed{
\sigma(U)
\subseteq
\{
z\in\mathbb{C}:|z|=1
\}.
}
\]

%%%%%%%%%%%%%%%%%%%%%%%%%%%%%%%%%%%%%%%%%%%%%%%%%%%%%%%%%%%%
\subsection{Isometries}
\label{subsec:isometries-operators}
%%%%%%%%%%%%%%%%%%%%%%%%%%%%%%%%%%%%%%%%%%%%%%%%%%%%%%%%%%%%

An operator \(T:H\to H\) is an isometry if

\[
\boxed{
\|Tx\|=\|x\|
}
\]

for every \(x\in H\).

Every unitary operator is an isometry.

However, an isometry need not be surjective.

Thus

\[
\boxed{
\text{unitary}
\Longrightarrow
\text{isometry},
}
\]

but the converse requires surjectivity.

%%%%%%%%%%%%%%%%%%%%%%%%%%%%%%%%%%%%%%%%%%%%%%%%%%%%%%%%%%%%
\subsection{Positive Operators}
\label{subsec:positive-operators}
%%%%%%%%%%%%%%%%%%%%%%%%%%%%%%%%%%%%%%%%%%%%%%%%%%%%%%%%%%%%

\begin{definitionbox}{Positive Operator}
A bounded operator \(T\) on a Hilbert space is positive if

\[
\boxed{
\langle Tx,x\rangle
\geq0
}
\]

for every \(x\in H\).
\end{definitionbox}

We write

\[
\boxed{
T\geq0.
}
\]

Every positive bounded operator is self-adjoint.

A basic example is

\[
\boxed{
T^*T\geq0.
}
\]

%%%%%%%%%%%%%%%%%%%%%%%%%%%%%%%%%%%%%%%%%%%%%%%%%%%%%%%%%%%%
\subsection{Positive Square Roots}
\label{subsec:positive-square-roots}
%%%%%%%%%%%%%%%%%%%%%%%%%%%%%%%%%%%%%%%%%%%%%%%%%%%%%%%%%%%%

If

\[
T\geq0,
\]

then there exists a unique positive operator \(S\) such that

\[
\boxed{
S^2=T.
}
\]

This operator is denoted

\[
\boxed{
T^{1/2}.
}
\]

Positive square roots are central to operator factorizations.

%%%%%%%%%%%%%%%%%%%%%%%%%%%%%%%%%%%%%%%%%%%%%%%%%%%%%%%%%%%%
\subsection{Absolute Value of an Operator}
\label{subsec:absolute-value-operator}
%%%%%%%%%%%%%%%%%%%%%%%%%%%%%%%%%%%%%%%%%%%%%%%%%%%%%%%%%%%%

For a bounded operator \(T\) on a Hilbert space, define

\[
\boxed{
|T|
=
(T^*T)^{1/2}.
}
\]

This is the operator analogue of the absolute value of a complex number.

It satisfies

\[
\boxed{
\|Tx\|
=
\||T|x\|.
}
\]

%%%%%%%%%%%%%%%%%%%%%%%%%%%%%%%%%%%%%%%%%%%%%%%%%%%%%%%%%%%%
\subsection{Polar Decomposition}
\label{subsec:polar-decomposition}
%%%%%%%%%%%%%%%%%%%%%%%%%%%%%%%%%%%%%%%%%%%%%%%%%%%%%%%%%%%%

The polar decomposition is an operator analogue of

\[
z
=
e^{i\theta}|z|.
\]

Under the standard Hilbert-space formulation, a bounded operator \(T\)
can be written

\[
\boxed{
T=U|T|,
}
\]

where \(U\) is a partial isometry with the appropriate initial and final
spaces.

If \(T\) is invertible, \(U\) is unitary.

%%%%%%%%%%%%%%%%%%%%%%%%%%%%%%%%%%%%%%%%%%%%%%%%%%%%%%%%%%%%
\subsection{Orthogonal Projections}
\label{subsec:orthogonal-projections}
%%%%%%%%%%%%%%%%%%%%%%%%%%%%%%%%%%%%%%%%%%%%%%%%%%%%%%%%%%%%

\begin{definitionbox}{Orthogonal Projection}
A bounded operator \(P:H\to H\) is an orthogonal projection if

\[
\boxed{
P^2=P
}
\]

and

\[
\boxed{
P^*=P.
}
\]
\end{definitionbox}

If \(M\subseteq H\) is a closed subspace, then every \(x\in H\) has a
unique decomposition

\[
\boxed{
x=Px+(I-P)x,
}
\]

where

\[
Px\in M
\]

and

\[
(I-P)x\in M^\perp.
\]

%%%%%%%%%%%%%%%%%%%%%%%%%%%%%%%%%%%%%%%%%%%%%%%%%%%%%%%%%%%%
\subsection{Projection Spectrum}
\label{subsec:projection-spectrum}
%%%%%%%%%%%%%%%%%%%%%%%%%%%%%%%%%%%%%%%%%%%%%%%%%%%%%%%%%%%%

If

\[
P^2=P,
\]

then

\[
P(P-I)=0.
\]

Therefore the polynomial

\[
p(z)=z(z-1)
\]

annihilates \(P\).

By spectral mapping,

\[
\boxed{
\sigma(P)\subseteq\{0,1\}.
}
\]

If \(P\neq0\) and \(P\neq I\), then typically

\[
\boxed{
\sigma(P)=\{0,1\}.
}
\]

%%%%%%%%%%%%%%%%%%%%%%%%%%%%%%%%%%%%%%%%%%%%%%%%%%%%%%%%%%%%
\subsection{Finite-Dimensional Spectral Theorem}
\label{subsec:finite-spectral-theorem}
%%%%%%%%%%%%%%%%%%%%%%%%%%%%%%%%%%%%%%%%%%%%%%%%%%%%%%%%%%%%

For a normal matrix \(A\), there exists a unitary matrix \(U\) and a
diagonal matrix \(D\) such that

\[
\boxed{
A=UDU^*.
}
\]

Equivalently, there exists an orthonormal basis consisting of
eigenvectors of \(A\).

The infinite-dimensional spectral theorem generalizes this idea, but
diagonal matrices are replaced by multiplication operators or spectral
measures.

%%%%%%%%%%%%%%%%%%%%%%%%%%%%%%%%%%%%%%%%%%%%%%%%%%%%%%%%%%%%
\subsection{Spectral Theorem for Compact Self-Adjoint Operators}
\label{subsec:compact-self-adjoint-spectral}
%%%%%%%%%%%%%%%%%%%%%%%%%%%%%%%%%%%%%%%%%%%%%%%%%%%%%%%%%%%%

\begin{theorem}[Spectral Theorem for Compact Self-Adjoint Operators]
Let \(T\) be compact and self-adjoint on a Hilbert space \(H\).

Then the nonzero spectrum consists of real eigenvalues of finite
multiplicity, with possible accumulation only at \(0\).

Moreover, \(H\) possesses an orthonormal basis adapted to the eigenspaces
of \(T\), together with the kernel when necessary.
\end{theorem}

Thus compact self-adjoint operators behave much like symmetric matrices.

%%%%%%%%%%%%%%%%%%%%%%%%%%%%%%%%%%%%%%%%%%%%%%%%%%%%%%%%%%%%
\subsection{Spectral Theorem for Bounded Self-Adjoint Operators}
\label{subsec:bounded-self-adjoint-spectral}
%%%%%%%%%%%%%%%%%%%%%%%%%%%%%%%%%%%%%%%%%%%%%%%%%%%%%%%%%%%%

The general spectral theorem cannot always be written as a sum over
eigenvectors.

Instead, a bounded self-adjoint operator admits a representation

\[
\boxed{
T
=
\int_{\sigma(T)}
\lambda\,dE(\lambda),
}
\]

where \(E\) is a projection-valued measure.

This is the continuous analogue of diagonalization.

%%%%%%%%%%%%%%%%%%%%%%%%%%%%%%%%%%%%%%%%%%%%%%%%%%%%%%%%%%%%
\subsection{Projection-Valued Measures}
\label{subsec:projection-valued-measures}
%%%%%%%%%%%%%%%%%%%%%%%%%%%%%%%%%%%%%%%%%%%%%%%%%%%%%%%%%%%%

A projection-valued measure assigns an orthogonal projection

\[
E(A)
\]

to suitable measurable subsets \(A\subseteq\R\).

The projections satisfy measure-like properties, including

\[
\boxed{
E(A\cap B)
=
E(A)E(B).
}
\]

The spectral theorem uses these projections to decompose a Hilbert space
according to spectral regions.

%%%%%%%%%%%%%%%%%%%%%%%%%%%%%%%%%%%%%%%%%%%%%%%%%%%%%%%%%%%%
\subsection{Spectral Functional Calculus}
\label{subsec:spectral-functional-calculus}
%%%%%%%%%%%%%%%%%%%%%%%%%%%%%%%%%%%%%%%%%%%%%%%%%%%%%%%%%%%%

If

\[
T
=
\int\lambda\,dE(\lambda),
\]

then suitable functions \(f\) may be applied to \(T\) through

\[
\boxed{
f(T)
=
\int
f(\lambda)\,dE(\lambda).
}
\]

For example,

\[
\boxed{
T^2
=
\int
\lambda^2\,dE(\lambda),
}
\]

and

\[
\boxed{
e^T
=
\int
e^\lambda\,dE(\lambda).
}
\]

This gives a rigorous meaning to functions of operators.

%%%%%%%%%%%%%%%%%%%%%%%%%%%%%%%%%%%%%%%%%%%%%%%%%%%%%%%%%%%%
\subsection{Multiplication Operators}
\label{subsec:multiplication-operators}
%%%%%%%%%%%%%%%%%%%%%%%%%%%%%%%%%%%%%%%%%%%%%%%%%%%%%%%%%%%%

Let

\[
m:X\to\mathbb{C}
\]

be measurable.

Define

\[
\boxed{
(M_mf)(x)
=
m(x)f(x).
}
\]

On \(L^2(X,\mu)\), if

\[
m\in L^\infty,
\]

then \(M_m\) is bounded and

\[
\boxed{
\|M_m\|
=
\|m\|_\infty.
}
\]

The spectrum is closely related to the essential range of \(m\):

\[
\boxed{
\sigma(M_m)
=
\operatorname{ess\,ran}(m).
}
\]

Multiplication operators are canonical models for normal operators.

%%%%%%%%%%%%%%%%%%%%%%%%%%%%%%%%%%%%%%%%%%%%%%%%%%%%%%%%%%%%
\subsection{Shift Operators}
\label{subsec:shift-operators}
%%%%%%%%%%%%%%%%%%%%%%%%%%%%%%%%%%%%%%%%%%%%%%%%%%%%%%%%%%%%

On

\[
\ell^2(\mathbb{N}),
\]

define the unilateral right shift by

\[
\boxed{
S(x_1,x_2,x_3,\ldots)
=
(0,x_1,x_2,x_3,\ldots).
}
\]

Then

\[
\boxed{
\|Sx\|_2=\|x\|_2,
}
\]

so \(S\) is an isometry.

However, \(S\) is not surjective because every vector in its range has
first coordinate \(0\).

Therefore \(S\) is not unitary.

This simple operator illustrates several infinite-dimensional phenomena
that cannot occur for square finite-dimensional isometries.

%%%%%%%%%%%%%%%%%%%%%%%%%%%%%%%%%%%%%%%%%%%%%%%%%%%%%%%%%%%%
\subsection{Spectrum of the Unilateral Shift}
\label{subsec:spectrum-unilateral-shift}
%%%%%%%%%%%%%%%%%%%%%%%%%%%%%%%%%%%%%%%%%%%%%%%%%%%%%%%%%%%%

The unilateral shift satisfies

\[
\boxed{
\|S\|=1.
}
\]

Its spectrum is

\[
\boxed{
\sigma(S)
=
\{
\lambda\in\mathbb{C}:|\lambda|\leq1
\}.
}
\]

However, \(S\) has no eigenvalues.

Thus the spectrum can be a large set even when the point spectrum is
empty.

%%%%%%%%%%%%%%%%%%%%%%%%%%%%%%%%%%%%%%%%%%%%%%%%%%%%%%%%%%%%
\subsection{Integral Operators}
\label{subsec:integral-operators}
%%%%%%%%%%%%%%%%%%%%%%%%%%%%%%%%%%%%%%%%%%%%%%%%%%%%%%%%%%%%

An integral operator has the form

\[
\boxed{
Tf(x)
=
\int_Y
K(x,y)f(y)\,d\mu(y),
}
\]

where \(K\) is called the kernel of the operator.

Integral operators appear naturally in

\[
\boxed{
\text{boundary-value problems},
\quad
\text{Green's functions},
\quad
\text{probability},
\quad
\text{quantum mechanics}.
}
\]

%%%%%%%%%%%%%%%%%%%%%%%%%%%%%%%%%%%%%%%%%%%%%%%%%%%%%%%%%%%%
\subsection{Hilbert--Schmidt Operators}
\label{subsec:hilbert-schmidt-operators}
%%%%%%%%%%%%%%%%%%%%%%%%%%%%%%%%%%%%%%%%%%%%%%%%%%%%%%%%%%%%

Suppose

\[
K\in L^2(X\times X).
\]

The associated integral operator

\[
Tf(x)
=
\int_XK(x,y)f(y)\,d\mu(y)
\]

is a Hilbert--Schmidt operator under standard hypotheses.

Its Hilbert--Schmidt norm is

\[
\boxed{
\|T\|_{\mathrm{HS}}
=
\left(
\int_{X\times X}
|K(x,y)|^2
\,d(\mu\times\mu)(x,y)
\right)^{1/2}.
}
\]

Hilbert--Schmidt operators are compact.

%%%%%%%%%%%%%%%%%%%%%%%%%%%%%%%%%%%%%%%%%%%%%%%%%%%%%%%%%%%%
\subsection{Compact Operators}
\label{subsec:compact-operators}
%%%%%%%%%%%%%%%%%%%%%%%%%%%%%%%%%%%%%%%%%%%%%%%%%%%%%%%%%%%%

\begin{definitionbox}{Compact Operator}
A bounded operator

\[
T:X\to Y
\]

is compact if it maps bounded subsets of \(X\) into relatively compact
subsets of \(Y\).
\end{definitionbox}

Equivalently, if \((x_n)\) is bounded in \(X\), then

\[
(Tx_n)
\]

has a convergent subsequence in \(Y\).

Finite-rank operators are compact.

%%%%%%%%%%%%%%%%%%%%%%%%%%%%%%%%%%%%%%%%%%%%%%%%%%%%%%%%%%%%
\subsection{Compact Operators as Infinite-Dimensional Matrices}
\label{subsec:compact-matrix-analogy}
%%%%%%%%%%%%%%%%%%%%%%%%%%%%%%%%%%%%%%%%%%%%%%%%%%%%%%%%%%%%

Compact operators behave in many respects like finite-dimensional
matrices.

For a compact operator \(T\) on an infinite-dimensional Banach space,
every nonzero spectral value is an eigenvalue.

Moreover, nonzero eigenvalues have finite algebraic multiplicity, and
the only possible accumulation point of the spectrum is

\[
\boxed{
0.
}
\]

This structure is much stronger than for arbitrary bounded operators.

%%%%%%%%%%%%%%%%%%%%%%%%%%%%%%%%%%%%%%%%%%%%%%%%%%%%%%%%%%%%
\subsection{Finite-Rank Approximation}
\label{subsec:finite-rank-approximation}
%%%%%%%%%%%%%%%%%%%%%%%%%%%%%%%%%%%%%%%%%%%%%%%%%%%%%%%%%%%%

A finite-rank operator has finite-dimensional range.

On Hilbert spaces, compact operators can be approximated in operator norm
by finite-rank operators.

Schematically,

\[
\boxed{
T
=
\lim_{n\to\infty}T_n,
}
\]

where each \(T_n\) has finite rank.

This explains why compact operators retain many matrix-like properties.

%%%%%%%%%%%%%%%%%%%%%%%%%%%%%%%%%%%%%%%%%%%%%%%%%%%%%%%%%%%%
\subsection{Singular Values}
\label{subsec:singular-values}
%%%%%%%%%%%%%%%%%%%%%%%%%%%%%%%%%%%%%%%%%%%%%%%%%%%%%%%%%%%%

For a compact operator \(T\) on a Hilbert space, the singular values are
the nonnegative eigenvalues of

\[
|T|
=
(T^*T)^{1/2}.
\]

They are commonly denoted

\[
\boxed{
s_1(T)\geq s_2(T)\geq\cdots\geq0.
}
\]

Singular values generalize matrix singular values to infinite dimensions.

They measure the strength of an operator along orthogonal directions.

%%%%%%%%%%%%%%%%%%%%%%%%%%%%%%%%%%%%%%%%%%%%%%%%%%%%%%%%%%%%
\subsection{Singular Value Decomposition}
\label{subsec:operator-svd}
%%%%%%%%%%%%%%%%%%%%%%%%%%%%%%%%%%%%%%%%%%%%%%%%%%%%%%%%%%%%

For suitable compact operators between Hilbert spaces, one obtains an
expansion

\[
\boxed{
Tx
=
\sum_{n}
s_n
\langle x,v_n\rangle u_n,
}
\]

where

\[
(u_n)
\]

and

\[
(v_n)
\]

are orthonormal systems and \(s_n\geq0\) are singular values.

This is the infinite-dimensional analogue of matrix singular value
decomposition.

%%%%%%%%%%%%%%%%%%%%%%%%%%%%%%%%%%%%%%%%%%%%%%%%%%%%%%%%%%%%
\subsection{Trace-Class Operators}
\label{subsec:trace-class-operators}
%%%%%%%%%%%%%%%%%%%%%%%%%%%%%%%%%%%%%%%%%%%%%%%%%%%%%%%%%%%%

A compact operator \(T\) is trace class if its singular values satisfy

\[
\boxed{
\sum_n s_n(T)<\infty.
}
\]

For a positive trace-class operator,

\[
\boxed{
\operatorname{tr}(T)
=
\sum_n
\langle Te_n,e_n\rangle,
}
\]

where \((e_n)\) is any orthonormal basis.

The trace is independent of the choice of orthonormal basis.

Trace-class operators are important in quantum mechanics and spectral
theory.

%%%%%%%%%%%%%%%%%%%%%%%%%%%%%%%%%%%%%%%%%%%%%%%%%%%%%%%%%%%%
\subsection{Fredholm Operators}
\label{subsec:fredholm-operators}
%%%%%%%%%%%%%%%%%%%%%%%%%%%%%%%%%%%%%%%%%%%%%%%%%%%%%%%%%%%%

\begin{definitionbox}{Fredholm Operator}
A bounded operator

\[
T:X\to Y
\]

between Banach spaces is Fredholm if

\[
\boxed{
\dim(\ker T)<\infty,
}
\]

the range of \(T\) is closed, and

\[
\boxed{
\dim(Y/\operatorname{Ran}(T))<\infty.
}
\]
\end{definitionbox}

The quotient

\[
Y/\operatorname{Ran}(T)
\]

measures the finite-dimensional failure of surjectivity.

%%%%%%%%%%%%%%%%%%%%%%%%%%%%%%%%%%%%%%%%%%%%%%%%%%%%%%%%%%%%
\subsection{Fredholm Index}
\label{subsec:fredholm-index}
%%%%%%%%%%%%%%%%%%%%%%%%%%%%%%%%%%%%%%%%%%%%%%%%%%%%%%%%%%%%

The index of a Fredholm operator is

\[
\boxed{
\operatorname{ind}(T)
=
\dim(\ker T)
-
\dim(Y/\operatorname{Ran}(T)).
}
\]

The index is stable under many perturbations.

In particular, compact perturbations preserve the Fredholm property and
index under standard hypotheses.

%%%%%%%%%%%%%%%%%%%%%%%%%%%%%%%%%%%%%%%%%%%%%%%%%%%%%%%%%%%%
\subsection{Fredholm Alternative}
\label{subsec:fredholm-alternative}
%%%%%%%%%%%%%%%%%%%%%%%%%%%%%%%%%%%%%%%%%%%%%%%%%%%%%%%%%%%%

For compact operators, equations of the form

\[
\boxed{
(I-K)x=y
}
\]

have a particularly useful solvability theory.

A standard Hilbert-space version of the Fredholm alternative says that,
for compact \(K\),

\[
I-K
\]

is invertible precisely when

\[
\boxed{
\ker(I-K)=\{0\}.
}
\]

More generally, solvability is controlled by the kernel of the adjoint
problem.

This is fundamental in integral equations and boundary-value problems.

%%%%%%%%%%%%%%%%%%%%%%%%%%%%%%%%%%%%%%%%%%%%%%%%%%%%%%%%%%%%
\subsection{Unbounded Operators}
\label{subsec:unbounded-operators}
%%%%%%%%%%%%%%%%%%%%%%%%%%%%%%%%%%%%%%%%%%%%%%%%%%%%%%%%%%%%

Many of the most important operators in analysis are not bounded.

Examples include

\[
\boxed{
Df=f',
}
\]

and

\[
\boxed{
-\Delta f.
}
\]

Such operators cannot generally be defined on the entire ambient Banach
or Hilbert space.

Instead, one specifies a domain

\[
\boxed{
\mathcal{D}(T)\subseteq X.
}
\]

Then

\[
\boxed{
T:\mathcal{D}(T)\to X.
}
\]

%%%%%%%%%%%%%%%%%%%%%%%%%%%%%%%%%%%%%%%%%%%%%%%%%%%%%%%%%%%%
\subsection{Why Domains Matter}
\label{subsec:operator-domains}
%%%%%%%%%%%%%%%%%%%%%%%%%%%%%%%%%%%%%%%%%%%%%%%%%%%%%%%%%%%%

For unbounded operators, the formula for \(T\) alone does not determine
the operator.

The domain is part of the definition.

For example,

\[
Tf=-f''
\]

with Dirichlet boundary conditions and

\[
Tf=-f''
\]

with periodic boundary conditions define different operators.

Thus

\[
\boxed{
\text{operator}
=
\text{formula}
+
\text{domain}.
}
\]

%%%%%%%%%%%%%%%%%%%%%%%%%%%%%%%%%%%%%%%%%%%%%%%%%%%%%%%%%%%%
\subsection{Densely Defined Operators}
\label{subsec:densely-defined-operators}
%%%%%%%%%%%%%%%%%%%%%%%%%%%%%%%%%%%%%%%%%%%%%%%%%%%%%%%%%%%%

An operator \(T\) on a Hilbert space \(H\) is densely defined if

\[
\boxed{
\overline{\mathcal{D}(T)}=H.
}
\]

Density is important because the adjoint of an unbounded operator is
defined naturally when the original operator is densely defined.

Differential operators on smooth compactly supported functions provide
standard examples.

%%%%%%%%%%%%%%%%%%%%%%%%%%%%%%%%%%%%%%%%%%%%%%%%%%%%%%%%%%%%
\subsection{Graph of an Operator}
\label{subsec:graph-operator}
%%%%%%%%%%%%%%%%%%%%%%%%%%%%%%%%%%%%%%%%%%%%%%%%%%%%%%%%%%%%

The graph of an operator

\[
T:\mathcal{D}(T)\subseteq X\to Y
\]

is

\[
\boxed{
G(T)
=
\{
(x,Tx):
x\in\mathcal{D}(T)
\}
\subseteq
X\times Y.
}
\]

Properties of an operator can often be studied geometrically through its
graph.

%%%%%%%%%%%%%%%%%%%%%%%%%%%%%%%%%%%%%%%%%%%%%%%%%%%%%%%%%%%%
\subsection{Closed Operators}
\label{subsec:closed-operators}
%%%%%%%%%%%%%%%%%%%%%%%%%%%%%%%%%%%%%%%%%%%%%%%%%%%%%%%%%%%%

\begin{definitionbox}{Closed Operator}
An operator \(T\) is closed if its graph

\[
G(T)
\]

is closed in

\[
X\times Y.
\]
\end{definitionbox}

Equivalently, if

\[
x_n\to x
\]

and

\[
Tx_n\to y,
\]

with \(x_n\in\mathcal{D}(T)\), then

\[
\boxed{
x\in\mathcal{D}(T)
\quad\text{and}\quad
Tx=y.
}
\]

%%%%%%%%%%%%%%%%%%%%%%%%%%%%%%%%%%%%%%%%%%%%%%%%%%%%%%%%%%%%
\subsection{Closed Graph Theorem}
\label{subsec:closed-graph-theorem-operator}
%%%%%%%%%%%%%%%%%%%%%%%%%%%%%%%%%%%%%%%%%%%%%%%%%%%%%%%%%%%%

\begin{theorem}[Closed Graph Theorem]
Let \(X\) and \(Y\) be Banach spaces.

If

\[
T:X\to Y
\]

is linear, defined on all of \(X\), and has closed graph, then \(T\) is
bounded.
\end{theorem}

The requirement that \(T\) be defined on all of \(X\) is essential.

Unbounded closed operators therefore necessarily have proper domains.

%%%%%%%%%%%%%%%%%%%%%%%%%%%%%%%%%%%%%%%%%%%%%%%%%%%%%%%%%%%%
\subsection{Closable Operators}
\label{subsec:closable-operators}
%%%%%%%%%%%%%%%%%%%%%%%%%%%%%%%%%%%%%%%%%%%%%%%%%%%%%%%%%%%%

An operator is closable if the closure of its graph is itself the graph
of an operator.

The resulting operator is called the closure and is denoted

\[
\boxed{
\overline{T}.
}
\]

Closability allows an initially defined operator to be extended to a
closed operator in a controlled manner.

%%%%%%%%%%%%%%%%%%%%%%%%%%%%%%%%%%%%%%%%%%%%%%%%%%%%%%%%%%%%
\subsection{Adjoints of Unbounded Operators}
\label{subsec:unbounded-adjoints}
%%%%%%%%%%%%%%%%%%%%%%%%%%%%%%%%%%%%%%%%%%%%%%%%%%%%%%%%%%%%

Let \(T\) be densely defined on a Hilbert space \(H\).

A vector \(y\in H\) belongs to \(\mathcal{D}(T^*)\) when there exists
\(z\in H\) such that

\[
\boxed{
\langle Tx,y\rangle
=
\langle x,z\rangle
}
\]

for every

\[
x\in\mathcal{D}(T).
\]

Then

\[
\boxed{
T^*y=z.
}
\]

Unlike the bounded case,

\[
\mathcal{D}(T^*)
\]

may be a proper subset of \(H\).

%%%%%%%%%%%%%%%%%%%%%%%%%%%%%%%%%%%%%%%%%%%%%%%%%%%%%%%%%%%%
\subsection{Symmetric and Self-Adjoint Unbounded Operators}
\label{subsec:symmetric-selfadjoint-unbounded}
%%%%%%%%%%%%%%%%%%%%%%%%%%%%%%%%%%%%%%%%%%%%%%%%%%%%%%%%%%%%

A densely defined operator \(T\) is symmetric if

\[
\boxed{
\langle Tx,y\rangle
=
\langle x,Ty\rangle
}
\]

for all

\[
x,y\in\mathcal{D}(T).
\]

Equivalently,

\[
\boxed{
T\subseteq T^*.
}
\]

Self-adjointness requires the stronger condition

\[
\boxed{
T=T^*,
}
\]

including equality of domains.

Thus

\[
\boxed{
\text{self-adjoint}
\Longrightarrow
\text{symmetric},
}
\]

but a symmetric unbounded operator need not be self-adjoint.

%%%%%%%%%%%%%%%%%%%%%%%%%%%%%%%%%%%%%%%%%%%%%%%%%%%%%%%%%%%%
\subsection{Differential Operators}
\label{subsec:differential-operators}
%%%%%%%%%%%%%%%%%%%%%%%%%%%%%%%%%%%%%%%%%%%%%%%%%%%%%%%%%%%%

Consider

\[
\boxed{
Tf=-f''.
}
\]

On an interval \([a,b]\), the spectral properties depend strongly on
boundary conditions.

For example, Dirichlet conditions impose

\[
\boxed{
f(a)=f(b)=0.
}
\]

The eigenvalue problem

\[
\boxed{
-f''=\lambda f
}
\]

then produces a discrete family of eigenvalues and eigenfunctions.

This connects operator theory with Sturm--Liouville theory.

%%%%%%%%%%%%%%%%%%%%%%%%%%%%%%%%%%%%%%%%%%%%%%%%%%%%%%%%%%%%
\subsection{Eigenfunction Expansions}
\label{subsec:eigenfunction-expansions}
%%%%%%%%%%%%%%%%%%%%%%%%%%%%%%%%%%%%%%%%%%%%%%%%%%%%%%%%%%%%

For suitable self-adjoint differential operators, eigenfunctions form an
orthogonal or orthonormal basis.

A function may then be expanded as

\[
\boxed{
f
=
\sum_n
\langle f,\phi_n\rangle\phi_n.
}
\]

This generalizes Fourier series.

Indeed, ordinary Fourier series arise from eigenfunctions of differential
operators such as

\[
-\frac{d^2}{dx^2}
\]

with suitable boundary conditions.

%%%%%%%%%%%%%%%%%%%%%%%%%%%%%%%%%%%%%%%%%%%%%%%%%%%%%%%%%%%%
\subsection{Operator Exponentials}
\label{subsec:operator-exponentials}
%%%%%%%%%%%%%%%%%%%%%%%%%%%%%%%%%%%%%%%%%%%%%%%%%%%%%%%%%%%%

For a bounded operator \(T\), define

\[
\boxed{
e^{tT}
=
\sum_{n=0}^{\infty}
\frac{t^nT^n}{n!}.
}
\]

The series converges in operator norm.

It satisfies

\[
\boxed{
e^{0T}=I
}
\]

and

\[
\boxed{
e^{(s+t)T}
=
e^{sT}e^{tT}.
}
\]

Thus operator exponentials naturally describe evolution equations.

%%%%%%%%%%%%%%%%%%%%%%%%%%%%%%%%%%%%%%%%%%%%%%%%%%%%%%%%%%%%
\subsection{Abstract Evolution Equations}
\label{subsec:abstract-evolution-equations}
%%%%%%%%%%%%%%%%%%%%%%%%%%%%%%%%%%%%%%%%%%%%%%%%%%%%%%%%%%%%

An ordinary differential equation in a Banach space may be written

\[
\boxed{
u'(t)=Au(t).
}
\]

If \(A\) is bounded, then

\[
\boxed{
u(t)=e^{tA}u(0).
}
\]

For unbounded operators, the theory becomes more subtle and leads to
strongly continuous semigroups.

%%%%%%%%%%%%%%%%%%%%%%%%%%%%%%%%%%%%%%%%%%%%%%%%%%%%%%%%%%%%
\subsection{Semigroups Preview}
\label{subsec:operator-semigroups-preview}
%%%%%%%%%%%%%%%%%%%%%%%%%%%%%%%%%%%%%%%%%%%%%%%%%%%%%%%%%%%%

A family

\[
(T(t))_{t\geq0}
\]

is a semigroup if

\[
\boxed{
T(0)=I
}
\]

and

\[
\boxed{
T(t+s)
=
T(t)T(s).
}
\]

A strongly continuous semigroup additionally satisfies

\[
\boxed{
T(t)x\to x
\qquad
\text{as }t\to0^+
}
\]

for every \(x\).

The generator of the semigroup is often an unbounded differential
operator.

Semigroup theory provides an operator-theoretic framework for evolution
PDEs.

%%%%%%%%%%%%%%%%%%%%%%%%%%%%%%%%%%%%%%%%%%%%%%%%%%%%%%%%%%%%
\subsection{Resolvents and Differential Equations}
\label{subsec:resolvent-differential-equations}
%%%%%%%%%%%%%%%%%%%%%%%%%%%%%%%%%%%%%%%%%%%%%%%%%%%%%%%%%%%%

Suppose one wishes to solve

\[
\boxed{
(T-\lambda I)x=y.
}
\]

If

\[
\lambda\in\rho(T),
\]

then

\[
\boxed{
x
=
R(\lambda,T)y.
}
\]

Thus the resolvent converts an operator equation into an inverse problem.

Resolvent estimates are fundamental in spectral theory and PDE analysis.

%%%%%%%%%%%%%%%%%%%%%%%%%%%%%%%%%%%%%%%%%%%%%%%%%%%%%%%%%%%%
\subsection{Operator Norm Convergence}
\label{subsec:operator-norm-convergence}
%%%%%%%%%%%%%%%%%%%%%%%%%%%%%%%%%%%%%%%%%%%%%%%%%%%%%%%%%%%%

A sequence \(T_n\in\mathcal{B}(X,Y)\) converges to \(T\) in operator norm
if

\[
\boxed{
\|T_n-T\|\to0.
}
\]

This means

\[
\boxed{
\sup_{\|x\|\leq1}
\|T_nx-Tx\|
\to0.
}
\]

Operator-norm convergence is a uniform form of convergence over the unit
ball.

%%%%%%%%%%%%%%%%%%%%%%%%%%%%%%%%%%%%%%%%%%%%%%%%%%%%%%%%%%%%
\subsection{Strong Operator Convergence}
\label{subsec:strong-operator-convergence}
%%%%%%%%%%%%%%%%%%%%%%%%%%%%%%%%%%%%%%%%%%%%%%%%%%%%%%%%%%%%

A sequence \(T_n\) converges strongly to \(T\) if

\[
\boxed{
T_nx\to Tx
}
\]

for every \(x\in X\).

This is written schematically as

\[
\boxed{
T_n\to T
\quad\text{strongly}.
}
\]

Operator-norm convergence implies strong convergence, but the converse
need not hold.

%%%%%%%%%%%%%%%%%%%%%%%%%%%%%%%%%%%%%%%%%%%%%%%%%%%%%%%%%%%%
\subsection{Weak Operator Convergence}
\label{subsec:weak-operator-convergence}
%%%%%%%%%%%%%%%%%%%%%%%%%%%%%%%%%%%%%%%%%%%%%%%%%%%%%%%%%%%%

On a Hilbert space, \(T_n\) converges to \(T\) in the weak operator
topology if

\[
\boxed{
\langle T_nx,y\rangle
\to
\langle Tx,y\rangle
}
\]

for every \(x,y\in H\).

Thus

\[
\boxed{
\text{norm convergence}
\Longrightarrow
\text{strong convergence}
\Longrightarrow
\text{weak operator convergence}.
}
\]

The reverse implications generally fail.

%%%%%%%%%%%%%%%%%%%%%%%%%%%%%%%%%%%%%%%%%%%%%%%%%%%%%%%%%%%%
\subsection{Worked Example: Strong but Not Norm Convergence}
\label{subsec:worked-strong-not-norm}
%%%%%%%%%%%%%%%%%%%%%%%%%%%%%%%%%%%%%%%%%%%%%%%%%%%%%%%%%%%%

\begin{workedexamplebox}{Strong Convergence Without Operator-Norm Convergence}

Let

\[
H=\ell^2(\mathbb{N}),
\]

and define

\[
P_n(x_1,x_2,\ldots)
=
(x_1,\ldots,x_n,0,0,\ldots).
\]

For every \(x\in\ell^2\),

\[
\|P_nx-x\|_2\to0.
\]

Therefore

\[
P_n\to I
\]

strongly.

However,

\[
\|I-P_n\|=1
\]

for every \(n\), since \(I-P_n\) acts as the identity on vectors supported
after the \(n\)th coordinate.

\begin{answerbox}
\[
\boxed{
P_n\to I
\text{ strongly, but }
P_n\not\to I
\text{ in operator norm}.
}
\]
\end{answerbox}
\end{workedexamplebox}

%%%%%%%%%%%%%%%%%%%%%%%%%%%%%%%%%%%%%%%%%%%%%%%%%%%%%%%%%%%%
\subsection{Numerical Approximation and Operators}
\label{subsec:numerical-operator-approximation}
%%%%%%%%%%%%%%%%%%%%%%%%%%%%%%%%%%%%%%%%%%%%%%%%%%%%%%%%%%%%

Many numerical methods replace an infinite-dimensional operator equation

\[
\boxed{
Tx=y
}
\]

with a finite-dimensional approximation

\[
\boxed{
T_nx_n=y_n.
}
\]

Examples include

\[
\boxed{
\text{finite differences},
\quad
\text{finite elements},
\quad
\text{spectral methods},
\quad
\text{Galerkin methods}.
}
\]

Operator theory helps determine whether the approximations are stable and
whether their spectra approximate the true spectrum.

%%%%%%%%%%%%%%%%%%%%%%%%%%%%%%%%%%%%%%%%%%%%%%%%%%%%%%%%%%%%
\subsection{Quantum-Mechanical Operators}
\label{subsec:quantum-operators}
%%%%%%%%%%%%%%%%%%%%%%%%%%%%%%%%%%%%%%%%%%%%%%%%%%%%%%%%%%%%

In mathematical formulations of quantum mechanics, observables are
represented by self-adjoint operators.

Examples include position,

\[
\boxed{
(Q\psi)(x)=x\psi(x),
}
\]

and momentum,

\[
\boxed{
(P\psi)(x)
=
-i\hbar\psi'(x),
}
\]

on suitable domains.

The symbol

\[
\hbar
\]

is read ``h-bar'' and denotes the reduced Planck constant.

The spectral theorem provides the mathematical framework for possible
measurement values.

%%%%%%%%%%%%%%%%%%%%%%%%%%%%%%%%%%%%%%%%%%%%%%%%%%%%%%%%%%%%
\subsection{Commutators}
\label{subsec:operator-commutators}
%%%%%%%%%%%%%%%%%%%%%%%%%%%%%%%%%%%%%%%%%%%%%%%%%%%%%%%%%%%%

For operators \(A\) and \(B\), define the commutator

\[
\boxed{
[A,B]
=
AB-BA.
}
\]

If

\[
[A,B]=0,
\]

then \(A\) and \(B\) commute.

Operator noncommutativity is fundamental in quantum mechanics.

For formal position and momentum operators,

\[
\boxed{
[Q,P]
=
i\hbar I.
}
\]

Domain issues must be handled carefully for unbounded operators.

%%%%%%%%%%%%%%%%%%%%%%%%%%%%%%%%%%%%%%%%%%%%%%%%%%%%%%%%%%%%
\subsection{Common Errors}
\label{subsec:operator-theory-common-errors}
%%%%%%%%%%%%%%%%%%%%%%%%%%%%%%%%%%%%%%%%%%%%%%%%%%%%%%%%%%%%

\subsubsection{Treating Every Operator Like a Finite Matrix}

Infinite-dimensional operators can have spectral values that are not
eigenvalues.

%%%%%%%%%%%%%%%%%%%%%%%%%%%%%%%%%%%%%%%%%%%%%%%%%%%%%%%%%%%%
\subsubsection{Assuming Every Linear Operator Is Bounded}

In infinite-dimensional normed spaces, linearity does not imply
boundedness.

Differentiation provides a standard example of an unbounded operator.

%%%%%%%%%%%%%%%%%%%%%%%%%%%%%%%%%%%%%%%%%%%%%%%%%%%%%%%%%%%%
\subsubsection{Ignoring the Domain of an Unbounded Operator}

For an unbounded operator,

\[
\boxed{
\text{formula alone does not define the operator}.
}
\]

The domain and boundary conditions are essential.

%%%%%%%%%%%%%%%%%%%%%%%%%%%%%%%%%%%%%%%%%%%%%%%%%%%%%%%%%%%%
\subsubsection{Confusing Symmetric and Self-Adjoint}

For unbounded operators,

\[
\boxed{
T\subseteq T^*
}
\]

does not imply

\[
T=T^*.
\]

Equality of domains must also hold.

%%%%%%%%%%%%%%%%%%%%%%%%%%%%%%%%%%%%%%%%%%%%%%%%%%%%%%%%%%%%
\subsubsection{Assuming Isometries Are Automatically Unitary}

An isometry preserves norms but may fail to be surjective.

The unilateral shift is the standard example.

%%%%%%%%%%%%%%%%%%%%%%%%%%%%%%%%%%%%%%%%%%%%%%%%%%%%%%%%%%%%
\subsubsection{Assuming the Spectrum Consists Only of Eigenvalues}

The point spectrum may be strictly smaller than the full spectrum.

%%%%%%%%%%%%%%%%%%%%%%%%%%%%%%%%%%%%%%%%%%%%%%%%%%%%%%%%%%%%
\subsubsection{Confusing Resolvent and Spectrum}

The resolvent set consists of values for which the shifted operator has a
bounded inverse.

The spectrum is its complement.

%%%%%%%%%%%%%%%%%%%%%%%%%%%%%%%%%%%%%%%%%%%%%%%%%%%%%%%%%%%%
\subsubsection{Assuming \(r(T)=\|T\|\) for Every Operator}

In general,

\[
\boxed{
r(T)\leq\|T\|.
}
\]

Equality holds for important classes such as normal operators on Hilbert
spaces, but not for arbitrary operators.

%%%%%%%%%%%%%%%%%%%%%%%%%%%%%%%%%%%%%%%%%%%%%%%%%%%%%%%%%%%%
\subsubsection{Assuming Strong Convergence Implies Norm Convergence}

Strong convergence is pointwise in vectors.

Operator-norm convergence is uniform over the unit ball and is strictly
stronger.

%%%%%%%%%%%%%%%%%%%%%%%%%%%%%%%%%%%%%%%%%%%%%%%%%%%%%%%%%%%%
\subsubsection{Ignoring Compactness Hypotheses}

Spectral results for compact operators do not automatically hold for
arbitrary bounded operators.

%%%%%%%%%%%%%%%%%%%%%%%%%%%%%%%%%%%%%%%%%%%%%%%%%%%%%%%%%%%%
\subsubsection{Using the Spectral Theorem Without Checking Operator Type}

Different versions of the spectral theorem apply to

\[
\boxed{
\text{normal},
\quad
\text{self-adjoint},
\quad
\text{unitary},
\quad
\text{compact self-adjoint}
}
\]

operators.

The hypotheses determine the correct conclusion.

%%%%%%%%%%%%%%%%%%%%%%%%%%%%%%%%%%%%%%%%%%%%%%%%%%%%%%%%%%%%
\subsection{Applications}
\label{subsec:operator-theory-applications}
%%%%%%%%%%%%%%%%%%%%%%%%%%%%%%%%%%%%%%%%%%%%%%%%%%%%%%%%%%%%

\begin{applicationbox}
\textbf{Differential Equations.}
Differential equations can be written as operator equations such as
\(Lu=f\), allowing invertibility and spectral methods to determine
existence and behavior of solutions.

\medskip

\textbf{Partial Differential Equations.}
Laplacians, elliptic operators, evolution operators, and semigroup
generators are naturally studied using operator theory.

\medskip

\textbf{Quantum Mechanics.}
Observables are represented by self-adjoint operators, while spectral
measures describe possible measurement values.

\medskip

\textbf{Integral Equations.}
Compact and Fredholm operators provide a framework for equations involving
integral kernels.

\medskip

\textbf{Fourier Analysis.}
Fourier transformation converts differential and convolution operators
into multiplication operators.

\medskip

\textbf{Harmonic Analysis.}
Singular integrals and Fourier multipliers are bounded operators on
function spaces.

\medskip

\textbf{Numerical Analysis.}
Finite-dimensional approximations of operators underlie finite-element,
spectral, and Galerkin methods.

\medskip

\textbf{Dynamical Systems.}
Transfer operators and evolution operators encode the long-term behavior
of dynamical systems.

\medskip

\textbf{Probability.}
Markov and transition operators describe evolution of probability
distributions.

\medskip

\textbf{Statistics and Data Science.}
Covariance operators and kernel integral operators extend matrix methods
to function spaces.
\end{applicationbox}

%%%%%%%%%%%%%%%%%%%%%%%%%%%%%%%%%%%%%%%%%%%%%%%%%%%%%%%%%%%%
\subsection{Exercises}
\label{subsec:operator-theory-exercises}
%%%%%%%%%%%%%%%%%%%%%%%%%%%%%%%%%%%%%%%%%%%%%%%%%%%%%%%%%%%%

\begin{exercise}[1]
Show that the identity operator satisfies

\[
\|I\|=1
\]

on every nonzero normed space.
\end{exercise}

\begin{exercise}[1]
Prove

\[
\|ST\|
\leq
\|S\|\|T\|.
\]
\end{exercise}

\begin{exercise}[1]
Show that

\[
\ker T
\]

is a linear subspace.
\end{exercise}

\begin{exercise}[1]
Show that

\[
\operatorname{Ran}(T)
\]

is a linear subspace.
\end{exercise}

\begin{exercise}[1]
Let

\[
Tf=2f
\]

on \(L^2([0,1])\). Compute \(\|T\|\) and \(\sigma(T)\).
\end{exercise}

\begin{exercise}[2]
Prove that boundedness and continuity are equivalent for linear
operators.
\end{exercise}

\begin{exercise}[2]
Prove the Neumann-series formula

\[
(I-T)^{-1}
=
\sum_{n=0}^{\infty}T^n
\]

when \(\|T\|<1\).
\end{exercise}

\begin{exercise}[2]
Show that

\[
\sigma(T)
\subseteq
\{
\lambda:|\lambda|\leq\|T\|
\}.
\]
\end{exercise}

\begin{exercise}[2]
Compute the spectrum of

\[
T=\lambda_0I.
\]
\end{exercise}

\begin{exercise}[2]
Show that every eigenvalue belongs to the spectrum.
\end{exercise}

\begin{exercise}[2]
Prove that a self-adjoint operator has real eigenvalues.
\end{exercise}

\begin{exercise}[2]
Show that eigenvectors of a self-adjoint operator corresponding to
distinct eigenvalues are orthogonal.
\end{exercise}

\begin{exercise}[2]
Prove that every unitary operator is an isometry.
\end{exercise}

\begin{exercise}[2]
Show that an orthogonal projection satisfies

\[
\|P\|\leq1.
\]

Determine when equality holds.
\end{exercise}

\begin{exercise}[2]
Show that

\[
\sigma(P)\subseteq\{0,1\}
\]

for a projection \(P\).
\end{exercise}

\begin{exercise}[2]
Verify that the unilateral shift on \(\ell^2\) is an isometry but is not
unitary.
\end{exercise}

\begin{exercise}[2]
Let

\[
(M_mf)(x)=m(x)f(x)
\]

on \(L^2\). Show that

\[
\|M_m\|=\|m\|_\infty.
\]
\end{exercise}

\begin{exercise}[3]
Prove that

\[
\|T^*T\|=\|T\|^2.
\]
\end{exercise}

\begin{exercise}[3]
Show that \(T^*T\) is positive.
\end{exercise}

\begin{exercise}[3]
Prove that every positive bounded operator is self-adjoint.
\end{exercise}

\begin{exercise}[3]
Study the construction of the positive square root

\[
T^{1/2}
\]

using the spectral theorem.
\end{exercise}

\begin{exercise}[3]
Prove the spectral radius formula for a bounded operator.
\end{exercise}

\begin{exercise}[3]
Prove the polynomial spectral mapping theorem.
\end{exercise}

\begin{exercise}[3]
Determine the point spectrum of the unilateral shift.
\end{exercise}

\begin{exercise}[3]
Show that a finite-rank operator is compact.
\end{exercise}

\begin{exercise}[3]
Show that the operator-norm limit of compact operators is compact.
\end{exercise}

\begin{exercise}[3]
Prove that a Hilbert--Schmidt integral operator is compact under standard
hypotheses.
\end{exercise}

\begin{exercise}[3]
Show that the coordinate projections

\[
P_n:\ell^2\to\ell^2
\]

converge strongly to \(I\) but not in operator norm.
\end{exercise}

\begin{exercise}[3]
For a bounded operator \(A\), show that

\[
u(t)=e^{tA}u_0
\]

solves

\[
u'(t)=Au(t).
\]
\end{exercise}

\begin{exercise}[3]
Show that

\[
e^{(s+t)A}
=
e^{sA}e^{tA}
\]

for a bounded operator \(A\).
\end{exercise}

\begin{exercise}[4]
Prove the spectral theorem for compact self-adjoint operators.
\end{exercise}

\begin{exercise}[4]
Study the projection-valued-measure form of the spectral theorem for
bounded self-adjoint operators.
\end{exercise}

\begin{exercise}[4]
Develop the continuous functional calculus for bounded self-adjoint
operators.
\end{exercise}

\begin{exercise}[4]
Prove that every nonzero spectral value of a compact operator is an
eigenvalue.
\end{exercise}

\begin{exercise}[4]
Study the Fredholm alternative for compact operators.
\end{exercise}

\begin{exercise}[4]
Show that the Fredholm index is stable under compact perturbations under
the standard hypotheses.
\end{exercise}

\begin{exercise}[4]
Study the polar decomposition theorem for bounded operators on Hilbert
spaces.
\end{exercise}

\begin{exercise}[4]
Investigate the spectral theorem for unbounded self-adjoint operators.
\end{exercise}

\begin{exercise}[4]
Study self-adjoint realizations of

\[
-\frac{d^2}{dx^2}
\]

under different boundary conditions.
\end{exercise}

\begin{exercise}[4]
Investigate the relationship between semigroup generators and evolution
equations.
\end{exercise}

%%%%%%%%%%%%%%%%%%%%%%%%%%%%%%%%%%%%%%%%%%%%%%%%%%%%%%%%%%%%
\subsection{Conceptual Summary}
\label{subsec:operator-theory-summary}
%%%%%%%%%%%%%%%%%%%%%%%%%%%%%%%%%%%%%%%%%%%%%%%%%%%%%%%%%%%%

Operator theory extends linear algebra to infinite-dimensional spaces.

A bounded linear operator satisfies

\[
\boxed{
\|Tx\|
\leq
\|T\|\|x\|.
}
\]

The spectrum is

\[
\boxed{
\sigma(T)
=
\{
\lambda\in\mathbb{C}:
T-\lambda I
\text{ is not boundedly invertible}
\}.
}
\]

Unlike finite-dimensional matrix theory,

\[
\boxed{
\text{spectrum}
\neq
\text{eigenvalues}
}
\]

in general.

The spectral radius is

\[
\boxed{
r(T)
=
\sup_{\lambda\in\sigma(T)}|\lambda|
=
\lim_{n\to\infty}\|T^n\|^{1/n}.
}
\]

On Hilbert spaces, the adjoint satisfies

\[
\boxed{
\langle Tx,y\rangle
=
\langle x,T^*y\rangle.
}
\]

Important classes include

\[
\boxed{
T=T^*
\quad
\text{self-adjoint},
}
\]

\[
\boxed{
TT^*=T^*T
\quad
\text{normal},
}
\]

and

\[
\boxed{
U^*U=UU^*=I
\quad
\text{unitary}.
}
\]

The spectral theorem generalizes diagonalization to infinite-dimensional
Hilbert spaces.

Compact operators behave particularly like finite-dimensional matrices.

Unbounded operators require explicit domains:

\[
\boxed{
T:\mathcal{D}(T)\subseteq H\to H.
}
\]

For these operators,

\[
\boxed{
\text{formula}
+
\text{domain}
}
\]

together determine the operator.

Operator theory therefore provides a common language for

\[
\boxed{
\text{linear algebra},
\quad
\text{functional analysis},
\quad
\text{differential equations},
\quad
\text{spectral theory},
\quad
\text{quantum mechanics}.
}
\]

%%%%%%%%%%%%%%%%%%%%%%%%%%%%%%%%%%%%%%%%%%%%%%%%%%%%%%%%%%%%
\subsection{Section Connections}
\label{subsec:operator-theory-connections}
%%%%%%%%%%%%%%%%%%%%%%%%%%%%%%%%%%%%%%%%%%%%%%%%%%%%%%%%%%%%

\chapterconnection{
Operator Theory develops the linear transformations introduced in Linear
Algebra within the infinite-dimensional framework of Functional Analysis.
The spectrum generalizes eigenvalue theory, while the spectral theorem
extends matrix diagonalization to Hilbert spaces. Fourier Analysis
provides important unitary and multiplication operators, and Harmonic
Analysis supplies singular integral and multiplier operators. Compact and
Fredholm operators lead naturally to integral equations and boundary-value
problems. Unbounded self-adjoint operators provide the mathematical
language for differential operators and quantum observables. The
operator-theoretic viewpoint is used extensively in Differential
Equations, Distribution Theory, Potential Theory, Geometric Analysis,
Numerical Analysis, Probability, and Mathematical Physics.
}

%%%%%%%%%%%%%%%%%%%%%%%%%%%%%%%%%%%%%%%%%%%%%%%%%%%%%%%%%%%%
% SECTION SOURCE TRACKING
%%%%%%%%%%%%%%%%%%%%%%%%%%%%%%%%%%%%%%%%%%%%%%%%%%%%%%%%%%%%

%------------------------------------------------------------
% 5.12 Operator Theory
%------------------------------------------------------------
%
% Source basis:
%   - Standard functional analysis and operator theory.
%   - Standard spectral theory on Banach and Hilbert spaces.
%   - Standard compact-operator and unbounded-operator theory.
%
% Used for:
%   - bounded linear operators;
%   - operator norms;
%   - operator algebras;
%   - kernels and ranges;
%   - bounded inverses;
%   - Neumann series;
%   - eigenvalues and eigenvectors;
%   - resolvent sets and resolvent operators;
%   - spectra;
%   - point, continuous, and residual spectrum;
%   - spectral radius;
%   - spectral radius formula;
%   - spectral mapping;
%   - functional calculus;
%   - adjoint operators;
%   - self-adjoint operators;
%   - normal operators;
%   - unitary operators;
%   - isometries;
%   - positive operators;
%   - positive square roots;
%   - operator absolute values;
%   - polar decomposition;
%   - orthogonal projections;
%   - spectral theorem;
%   - projection-valued measures;
%   - multiplication operators;
%   - shift operators;
%   - integral operators;
%   - Hilbert--Schmidt operators;
%   - compact operators;
%   - singular values;
%   - operator SVD;
%   - trace-class operators;
%   - Fredholm operators;
%   - Fredholm index;
%   - Fredholm alternative;
%   - unbounded operators;
%   - operator domains;
%   - closed and closable operators;
%   - unbounded adjoints;
%   - symmetric and self-adjoint operators;
%   - differential operators;
%   - eigenfunction expansions;
%   - operator exponentials;
%   - semigroup preview;
%   - operator convergence;
%   - numerical approximation;
%   - quantum-mechanical operators;
%   - commutators.
%
% Notes:
%   - Functional Analysis supplies the Banach-space and Hilbert-space
%     framework assumed throughout this section.
%   - Fourier Analysis provides fundamental examples of unitary
%     transformations and multiplication operators.
%   - Harmonic Analysis provides singular integral and Fourier
%     multiplier examples.
%   - Differential Equations later use unbounded operators,
%     eigenfunction expansions, and semigroups extensively.
%   - Distribution Theory provides natural domains for generalized
%     differential operators.
%   - Spectral methods reappear throughout Mathematical Physics and
%     Numerical Analysis.
%
%%%%%%%%%%%%%%%%%%%%%%%%%%%%%%%%%%%%%%%%%%%%%%%%%%%%%%%%%%%%